% This is an annotated latex definition file.
% You can give it any name as long as its inputted somewhere
% in your code.
% You may even put it in a directory other then the current
% as long as that directory is in the TEXINPUTS path.
% Comments start with a % as you a seeing right now. Otherwise look
% to the left.
% font of the times family
%\usepackage{palatino}
\usepackage{lmodern}
% make linenumbering possible for reviewing
\usepackage[pagewise,mathlines,displaymath]{lineno}
% use the pdftex graphics extensions
\usepackage{graphicx}
\usepackage{wrapfig}
% to make the headers and footers look nice
\usepackage{xcolor,fancyhdr}
% allow drawing with tikz
\usepackage{tikz}
% allow compact lists etc.
\usepackage{enumitem}
% make the header (and footer) span the whole printable area
\addtolength\headwidth\marginparwidth
\addtolength\headwidth\marginparsep
\pagestyle{fancy}
\usepackage[all]{svn-multi}
\svnRegisterAuthor{879417}{Pieter van den Hombergh}

% \usepackage{svn}
\newcommand\TheFile{mydefs.tex}
\newcommand\Code[1]{\textbf{\texttt{#1}\ }}
%%
% to get version information in files, include the following lines in
% each document 
% and renewcommand\TheFile to something like \renewcommand\TheFile{main.text}
%\renewcommand\footrulewidth{1pt}
% put revision control info in the footer
% the vspace compensates for the space taken up by the logo on the right hand.
% I like a fontys color in my header
\definecolor{purple}{rgb}{0.21,0,0.21} % a darkish purple
\renewcommand\headrule{%
  {\color{purple}%
    \hrule height 2pt width \headwidth
    \vspace{1pt}%
    \hrule height 1pt width \headwidth
    \vspace{-4pt}%
  }%
}%
\renewcommand\footrule{
  {\color{purple}%
    \hrule height 1pt width \headwidth
    \vspace{1pt}%
    \hrule height 2pt width \headwidth
    \vspace{2pt}%
  }%
}
% do not indent pars, increase parskip
\setlength\parindent{0pt}
\setlength\parskip{0.5em}
% my own itemize: a bit less spacy then the default
\newenvironment{Itemize} {
  \begin{itemize}{}%
    \setlength\topsep{0ex}%
    \setlength\parskip{0ex}%
    \setlength\partopsep{0em}%
    \setlength\parsep{0em}%
    \setlength\itemsep{0em}%
    }%
  {\end{itemize}}
% make the head a bit heigher
\setlength\headheight{16pt}
\setlength\footskip{60pt}
%%
% This macro '\define' puts the argument in em 
% and in boldface in the margin.
\newcommand{\define}[1]{% 1 argument
  \mbox{}{\textit{#1}}% italics or em
  \marginpar{\raggedright% no adjust
    \bfseries\hspace{0pt}#1}% bold
} % end of macro

%%
% At start of pages, define sometimes puts the material in the wrong margin
% in that case use this fix.
\newcommand{\fixdefine}[1]{\mbox{}{\textit{#1}}%
  \reversemarginpar\marginpar{\bfseries\hspace{0pt}#1}\normalmarginpar}
% to display source code like things
\usepackage{listings}
\lstset{numbers=right} % to get out of the way with document line numbers
% to be able to 'on this page' and the like
\usepackage{varioref}
% to put a nice quote with the chapter
\usepackage[avantgarde]{quotchap}
\renewcommand\chapterheadstartvskip{\vspace*{-5\baselineskip}}
%select Helvetica for title and quote
\usepackage{helvet}
\renewcommand\sectfont{\sffamily\bfseries}

%% example
%\begin{savequote}[10cm]
% \sffamily
%some text
%\end{savequote}
% needed for inclusion of gnumeric table
\def\inputGnumericTable{}
%% 
\usepackage{array}    
\usepackage{longtable}
\usepackage{calc}     
\usepackage{multirow} 
\usepackage{hhline}   
\usepackage{ifthen}   
%%
% Fancy headers are used
% To make the footer appear on chapter pages too,
% we redefine the fancypagestyle plain
\fancypagestyle{plain}{%
  \fancyhf{} % clear all  
  \fancyfoot[RO,LE]{\vspace{-1.2cm}%
  File: \svnkw{Filename}\\%
  Author: \svnFullAuthor*{\svnfileauthor}\\%
  Revision: \svnrev, \svnyear-\svnmonth-\svnday\ \svnhour:\svnminute\\
}
  \fancyfoot[RE,LO]{\includegraphics[width=2cm]{figures/fon000_00c.pdf}}
  \fancyfoot[C]{\bfseries\thepage}
  \renewcommand\headrule{}%
}% end of redef fancypagestyl plain
% Define my own fancy headers and footers  
\fancyhead{}
\fancyhead[RO]{\rightmark} % display section name
\fancyhead[LE]{\rightmark}
\fancyfoot[RO,LE]{\vspace{-1.2cm}%
  File: \svnkw{Filename}\\%
  Author: \svnFullAuthor*{\svnfileauthor}\\%
  Revision: \svnrev, \svnyear-\svnmonth-\svnday\ \svnhour:\svnminute\\
}
% logo in footer
\fancyfoot[RE,LO]{\includegraphics[width=2cm]{figures/fon000_00c.pdf}}
\pagestyle{fancy}
\usepackage{fontyscoverpage}
%%
% For a watermark on each page:
\usepackage{eso-pic} % needed package
% Draw a DRAFT through the text.
% for other text renewcommand the next command
\newcommand\watermarktext{Draft}
% You could make this depend on the status of the files.
\makeatletter
  \AddToShipoutPicture{%
    \setlength{\@tempdimb}{.5\paperwidth}%
    \setlength{\@tempdimc}{.5\paperheight}%
    \setlength{\unitlength}{1pt}%
    \put(\strip@pt\@tempdimb,\strip@pt\@tempdimc){%
      \makebox(0,0){\rotatebox{55}{% for longer texts you might want to increase this angle
          \textcolor[gray]{0.85}{% the higher the number, the lighter the text
            \fontsize{8cm}{8cm}% Fit this font to the text
            \sffamily
            \selectfont{\watermarktext}}}}
    }
}
\makeatother
\lstdefinelanguage{BibTeX}
  {keywords={%
      @article,@book,@collectedbook,@conference,@electronic,@ieeetranbstctl,%
      @inbook,@incollectedbook,@incollection,@injournal,@inproceedings,%
      @manual,@mastersthesis,@misc,@patent,@periodical,@phdthesis,@preamble,%
      @proceedings,@standard,@string,@techreport,@unpublished%
      },
   comment=[l][\itshape]{@comment},
   sensitive=false,
  }

% thenext comments are inserted by Emacs, that uses it 
% to process this latex file
%%% Local Variables: 
%%% mode: latex
%%% TeX-master: t
%%% End: 
